\section*{Capítulo \ref{ch:b1:matrices} --- Matrices}

\begin{solution}{exo:b1:matrices:1}
\[
AB = \begin{pmatrix} 2 & 1\\ 1 & 0\end{pmatrix},
\quad
BA = \begin{pmatrix} 0 & 1\\ 1 & 2\end{pmatrix},
\quad
A^2 - B^2 = \begin{pmatrix} 1 & 4\\ 0 & 1\end{pmatrix} - I
= \begin{pmatrix} 0 & 4\\ 0 & 0\end{pmatrix},
\]
\[
(A+B)(A-B) = A^2 - AB + BA - B^2
= \begin{pmatrix} 0 & 4\\ 0 & 0\end{pmatrix} +
\begin{pmatrix} -2 & 0\\ 0 & 2 \end{pmatrix}
= \begin{pmatrix} -2 & 4\\ 0 & 2\end{pmatrix}.
\]
Difieren en $BA - AB \neq 0$: la identidad
$(a+b)(a-b) = a^2 - b^2$ exige la conmutatividad, que aquí falla.
\end{solution}

\begin{solution}{exo:b1:matrices:2}
$M$: $L_2 \leftarrow L_2 - 2L_1$ mata la segunda fila;
$L_3 \leftarrow L_3 - L_1$ da $(0, -1, -2)$. Dos pivotes:
$\operatorname{rk} M = 2$.

$N$: $L_3 \leftarrow L_3 - L_1$ da $(0,1,1,1) = L_2$; después
$L_3 \leftarrow L_3 - L_2 = 0$. Dos pivotes: $\operatorname{rk} N = 2$.
\end{solution}

\begin{solution}{exo:b1:matrices:3}
Reduciendo $(A \mid I_3)$: $L_2 \leftarrow L_2 - 2L_1$,
$L_3 \leftarrow L_3 - L_1$:
\[
\begin{pmatrix}
1 & 0 & 1 & 1 & 0 & 0\\
0 & 1 & -1 & -2 & 1 & 0\\
0 & 1 & 0 & -1 & 0 & 1
\end{pmatrix}
\xrightarrow{L_3 \leftarrow L_3 - L_2}
\begin{pmatrix}
1 & 0 & 1 & 1 & 0 & 0\\
0 & 1 & -1 & -2 & 1 & 0\\
0 & 0 & 1 & 1 & -1 & 1
\end{pmatrix},
\]
después $L_1 \leftarrow L_1 - L_3$, $L_2 \leftarrow L_2 + L_3$:
\[
A^{-1} = \begin{pmatrix}
0 & 1 & -1\\
-1 & 0 & 1\\
1 & -1 & 1
\end{pmatrix}.
\]
\emph{Comprobación:} la primera fila de $A$ por la primera columna de
$A^{-1}$: $1 \cdot 0 + 0\cdot(-1) + 1\cdot 1 = 1$; por la segunda
columna: $1 - 0 - 1 = 0$; por la tercera: $-1 + 0 + 1 = 0$.
\end{solution}

\begin{solution}{exo:b1:matrices:4}
$u(1) = 1$, $u(X) = X + 1$, $u(X^2) = X^2 + 2X + 1$: las columnas de
\omterm{prop:b1:vspaces:coordinates}{coordenadas} en $(1, X, X^2)$ dan
\[
M = \begin{pmatrix}
1 & 1 & 1\\
0 & 1 & 2\\
0 & 0 & 1
\end{pmatrix}.
\]
$u$ es invertible porque tiene la inversa obvia $P \mapsto P(X - 1)$
(composición de sustituciones). Su matriz se obtiene igual a partir
de $u^{-1}(X^k) = (X-1)^k$:
\[
M^{-1} = \begin{pmatrix}
1 & -1 & 1\\
0 & 1 & -2\\
0 & 0 & 1
\end{pmatrix}.
\]
\end{solution}

\begin{solution}{exo:b1:matrices:5}
$N = \begin{pmatrix} 0 & 1\\ 0 & 0\end{pmatrix}$, $N^2 = 0$. Como
$2I$ y $N$ conmutan, el desarrollo binomial se trunca tras dos
términos:
\[
A^k = (2I + N)^k = 2^k I + k\,2^{k-1} N
= \begin{pmatrix} 2^k & k\,2^{k-1}\\ 0 & 2^k\end{pmatrix}.
\]
(Compruébese $k = 2$:
$A^2 = \begin{pmatrix}4 & 4\\ 0 & 4\end{pmatrix}$, correcto por
producto directo.)
\end{solution}

\begin{solution}{exo:b1:matrices:6}
\omterm{def:b1:matrices:transpose}{Trazas}: $\operatorname{tr}(AB - BA) = \operatorname{tr}(AB) -
\operatorname{tr}(BA) = 0$ (\cref{def:b1:matrices:transpose}),
mientras que $\operatorname{tr}(I_n) = n \neq 0$ en $\R$ o en $\C$.
No hay solución. (En espacios de dimensión infinita, la identidad
\emph{sí} es realizable ---la derivación y la multiplicación por $x$
la cumplen--- precisamente porque allí no existe ninguna \omterm{def:b1:matrices:transpose}{traza}.)
\end{solution}

\begin{solution}{exo:b1:matrices:7}
Producto telescópico, conmutando todas las potencias de $A$:
\[
(I - A)(I + A + \dots + A^{m-1}) = I - A^m = I ,
\]
y la~\cref{prop:b1:matrices:ring} eleva la inversa por un lado. Para
la \omterm{def:b1:logic:map}{aplicación}: la matriz dada es $I + N$ con
\[
N = \begin{pmatrix} 0 & 2 & 3\\ 0 & 0 & 2\\ 0&0&0 \end{pmatrix},
\quad
N^2 = \begin{pmatrix} 0&0&4\\ 0&0&0\\ 0&0&0\end{pmatrix},
\quad N^3 = 0 ,
\]
luego, sustituyendo $A$ por $-N$ en la fórmula:
\[
(I + N)^{-1} = I - N + N^2 =
\begin{pmatrix}
1 & -2 & 1\\
0 & 1 & -2\\
0 & 0 & 1
\end{pmatrix}.
\]
\end{solution}

\begin{solution}{exo:b1:matrices:8}
$A^2 = A$: el endomorfismo $a$ es una \omterm{def:b1:linmaps:projection}{proyección}
(\cref{thm:b1:linmaps:projchar}),
$E = \operatorname{im} a \oplus \ker a$ con
$\dim\operatorname{im} a = r = \operatorname{rk} A$. En una \omterm{def:b1:vspaces:free}{base}
adaptada a esa descomposición ($r$ vectores de la \omterm{def:b1:linmaps:kerim}{imagen} y después
una \omterm{def:b1:vspaces:free}{base} del núcleo), la matriz de $a$ es
$\begin{pmatrix} I_r & 0\\ 0 & 0\end{pmatrix}$, de \omterm{def:b1:matrices:transpose}{traza} $r$. Y la
\omterm{def:b1:matrices:transpose}{traza} es invariante por \omterm{def:b1:matrices:changeofbasis}{cambio de base}:
$\operatorname{tr}(P^{-1}MP) = \operatorname{tr}(MPP^{-1}) =
\operatorname{tr} M$ por la identidad cíclica. Por tanto,
$\operatorname{tr} A = r = \operatorname{rk} A$.
\end{solution}

\begin{solution}{exo:b1:matrices:9}
$J^2 = nJ$ (cada entrada de $J^2$ suma $n$ unos). Búsquese
$M^{-1} = \alpha I + \beta J$:
\[
(aI + bJ)(\alpha I + \beta J)
= a\alpha\, I + (a\beta + b\alpha + nb\beta)\, J .
\]
Esto es igual a $I$ si y solo si $a\alpha = 1$ y
$a\beta + b\alpha + nb\beta = 0$, es decir, $\alpha = \frac1a$ y
$\beta(a + nb) = -\frac ba$. Si $a \neq 0$ y $a + nb \neq 0$:
\[
M^{-1} = \frac 1a I - \frac{b}{a(a + nb)}\, J .
\]
Recíprocamente, si $a = 0$: $M = bJ$ tiene rango $\leq 1 < n$ (para
$n \geq 2$): no invertible ($n = 1$ es el caso escalar). Si
$a + nb = 0$: el vector $v = (1, \dots, 1)^{\mathsf T}$ cumple
$Mv = (a + nb)v = 0$ con $v \neq 0$: no invertible. Luego
$M \in GL_n \iff a \neq 0$ y $a + nb \neq 0$.
\end{solution}

\begin{solution}{exo:b1:matrices:10}
\emph{Suma:}
$\operatorname{im}(A + B) \subseteq \operatorname{im} A +
\operatorname{im} B$ (cada $(A+B)x = Ax + Bx$), y Grassmann acota la
dimensión de una suma por la suma de las dimensiones.

\emph{Sylvester:} sea $a$ la \omterm{def:b1:logic:map}{aplicación} de $A$ restringida a
$V = \operatorname{im} B$ (de dimensión $\operatorname{rk} B$). Su
\omterm{def:b1:linmaps:kerim}{imagen} es $\operatorname{im}(AB)$ ($a(Bx) = ABx$), y el teorema del
rango en $V$:
\[
\operatorname{rk} B = \dim\ker(a_{|V}) + \operatorname{rk}(AB) .
\]
Ahora bien, $\ker(a_{|V}) \subseteq \ker A$, de dimensión
$n - \operatorname{rk} A$: luego
\[
\operatorname{rk}(AB) \geq \operatorname{rk} B - (n -
\operatorname{rk} A) = \operatorname{rk} A + \operatorname{rk} B -
n . \qedhere
\]
\end{solution}

\begin{solution}{exo:b1:matrices:11}
\begin{enumerate}
  \item Entrada a entrada, $(AD)_{ij} = a_{ij}\,d_j$ y
        $(DA)_{ij} = d_i\,a_{ij}$. Luego $AD = DA$ si y solo si
        $a_{ij}(d_j - d_i) = 0$ para todos $i, j$; y cuando
        $i \neq j$ el factor $d_j - d_i$ es no nulo, lo que fuerza
        $a_{ij} = 0$: $A$ es diagonal. Recíprocamente, las matrices
        diagonales conmutan entre sí.
  \item Si $A$ conmuta con todas las matrices, conmuta con
        $\operatorname{diag}(1, 2, \dots, n)$, luego
        $A = \operatorname{diag}(\lambda_1, \dots, \lambda_n)$ por
        (1). Entonces $A E_{ij} = \lambda_i E_{ij}$ (solo sobrevive
        la fila $i$ de $E_{ij}$), mientras que
        $E_{ij} A = \lambda_j E_{ij}$: conmutar con $E_{ij}$ fuerza
        $\lambda_i = \lambda_j$. Por tanto, $A = \lambda I_n$; y las
        matrices escalares sí conmutan con todo. El centro de
        $\mathcal{M}_n(K)$ es $K\,I_n$.
\end{enumerate}
\end{solution}

\begin{solution}{exo:b1:matrices:12}
\begin{enumerate}
  \item Si $\operatorname{rk} A = 1$: la \omterm{def:b1:linmaps:kerim}{imagen} de $A$ es una recta
        $\operatorname{Vect}(C)$, $C \neq 0$, de modo que la columna
        $j$-ésima de $A$ es $\ell_j\,C$ para ciertos escalares
        $\ell_j$ (no todos nulos), es decir, $A = C L$ con
        $L = (\ell_1, \dots, \ell_n) \neq 0$. Recíprocamente, si
        $A = CL \neq 0$, todas las columnas son múltiplos de $C$:
        rango $1$.
  \item $A^2 = C\,(L C)\,L$, y $LC$ es el escalar
        $\sum_i \ell_i c_i = \operatorname{tr}(CL) =
        \operatorname{tr} A$. Luego $A^2 = (\operatorname{tr} A)\,A$
        y, por inducción, $A^m = (\operatorname{tr} A)^{m-1} A$. Si
        $\operatorname{tr} A \neq 0$, ninguna potencia se anula; y
        si $\operatorname{tr} A = 0$, entonces $A^2 = 0$: una matriz
        de rango uno es nilpotente si y solo si su \omterm{def:b1:matrices:transpose}{traza} es nula.
  \item Con $t = \operatorname{tr} A \neq -1$:
        \[
        (I_n + A)\Bigl(I_n - \frac{A}{1 + t}\Bigr)
        = I_n + A - \frac{A + A^2}{1 + t}
        = I_n + A - \frac{(1 + t)A}{1 + t} = I_n ,
        \]
        usando $A^2 = tA$. Si $t = -1$:
        $(I_n + A)A = A + A^2 = A - A = 0$ con $A \neq 0$, de modo
        que $I_n + A$ mata todas las columnas (no nulas) de $A$: no
        es \omterm{def:b1:logic:inj}{inyectiva} ni invertible.
\end{enumerate}
\end{solution}

\begin{solution}{pb:b1:matrices:1}
\textbf{1.} División euclídea de $X^n$ entre el \omterm{def:b1:poly:def}{polinomio mónico}
de grado $2$, $D$ (\cref{thm:b1:poly:division}): el cociente y el resto
existen y son únicos, y el resto tiene grado $\leq 1$:
$X^n = Q_n D + a_n X + b_n$. Para $n = 0$: $Q_0 = 0$,
$(a_0, b_0) = (0, 1)$; y para $n = 1$: $(a_1, b_1) = (1, 0)$.

\textbf{2.} Multiplíquese por $X$ y redúzcase $X^2 = D + sX - p$:
\[
X^{n+1} = X Q_n D + a_n X^2 + b_n X
= (X Q_n + a_n)\,D + (s\,a_n + b_n)\,X - p\,a_n .
\]
La última expresión tiene forma de resto (grado $\leq 1$), así que,
por unicidad, $a_{n+1} = s a_n + b_n$ y $b_{n+1} = -p a_n$.
Sustituyendo $b_{n+1} = -pa_n$ en $a_{n+2} = s a_{n+1} + b_{n+1}$ se
obtiene $a_{n+2} = s\,a_{n+1} - p\,a_n$.

\textbf{3.} Evalúese $X^n = Q_n D + a_n X + b_n$ en las raíces:
$\lambda^n = a_n\lambda + b_n$ y $\mu^n = a_n\mu + b_n$. Restando y
dividiendo por $\lambda - \mu \neq 0$:
\[
a_n = \frac{\lambda^n - \mu^n}{\lambda - \mu},
\qquad
b_n = \lambda^n - a_n\lambda
= \frac{\lambda\mu^n - \mu\lambda^n}{\lambda - \mu} .
\]

\textbf{4.} En la raíz doble: $\lambda^n = a_n\lambda + b_n$.
Derivando la identidad,
$nX^{n-1} = Q_n'\,(X - \lambda)^2 + 2Q_n\,(X - \lambda) + a_n$ y
evaluando en $\lambda$: $a_n = n\lambda^{n-1}$; y entonces
$b_n = \lambda^n - n\lambda^{n} = (1 - n)\lambda^{n}$.

\textbf{5.} Para $P = \sum_i p_i X^i$ y $Q = \sum_j q_j X^j$,
\[
P(M)\,Q(M) = \sum_{i,j} p_i q_j M^{i+j} = (PQ)(M),
\]
porque las potencias de una misma matriz $M$ conmutan entre sí (las
sumas son claras por linealidad). Si $D(M) = 0$, sustituir $M$ en
$X^n = Q_n D + a_n X + b_n$ da
$M^n = Q_n(M)\,D(M) + a_n M + b_n I = a_n M + b_n I$.

\textbf{6.} Productos directos:
\[
A^2 = \begin{pmatrix}
a^2 + bc & b(a + d)\\
c(a + d) & d^2 + bc
\end{pmatrix},
\qquad
s A = \begin{pmatrix}
a(a+d) & b(a+d)\\
c(a+d) & d(a+d)
\end{pmatrix},
\]
de modo que $A^2 - sA$ tiene nulas las entradas de fuera de la
diagonal y, en la diagonal, $a^2 + bc - a^2 - ad = bc - ad = -p$:
$A^2 - sA + pI_2 = 0$.

\textbf{7.} Con
$A' = \begin{pmatrix} a' & b'\\ c' & d'\end{pmatrix}$, desarrollando
$p(AA') = (aa' + bc')(cb' + dd') - (ab' + bd')(ca' + dc')$: los
términos $aa'cb'$ y $ab'ca'$ se cancelan, los términos $bc'dd'$ y
$bd'dc'$ se cancelan, y lo que queda es
\[
aa'dd' - bca'd' + bcb'c' - adb'c'
= (ad - bc)(a'd' - b'c') = p(A)\,p(A').
\]
Si $p \neq 0$, Cayley--Hamilton da
$A\,\bigl(\tfrac1p(sI_2 - A)\bigr) = \tfrac1p(sA - A^2) = I_2$, de
donde la inversa (y la~\cref{prop:b1:matrices:ring} la vuelve
bilátera). Y si $p = 0$ y $A$ fuera invertible, la multiplicatividad
daría $1 = p(I_2) = p(A)\,p(A^{-1}) = 0$: imposible. Luego
$A \in GL_2 \iff p \neq 0$.

\textbf{8.} $s = 3$, $p = 2$, $D = X^2 - 3X + 2 = (X - 1)(X - 2)$:
$\lambda = 2$, $\mu = 1$, luego $a_n = 2^n - 1$ y $b_n = 2 - 2^n$
(pregunta 3). Por tanto,
\[
A^n = (2^n - 1)A + (2 - 2^n)I
= \begin{pmatrix} 1 & 2^n - 1\\ 0 & 2^n \end{pmatrix}.
\]
Comprobación: $A^2 = \begin{pmatrix} 1 & 3\\ 0 & 4\end{pmatrix}$,
tanto por la fórmula como elevando al cuadrado directamente.

\textbf{9.} $s = 4$, $p = 3\cdot1 - 1\cdot(-1) = 4$:
$D = X^2 - 4X + 4 = (X - 2)^2$, con raíz doble $\lambda = 2$.
Pregunta 4: $a_n = n\,2^{n-1}$, $b_n = (1 - n)2^n$, luego
\[
A^n = n\,2^{n-1}A + (1 - n)2^n I
= 2^{n-1}\begin{pmatrix} n + 2 & n\\ -n & 2 - n
\end{pmatrix}.
\]
En $n = 2$:
$2\begin{pmatrix} 4 & 2\\ -2 & 0\end{pmatrix} =
\begin{pmatrix} 8 & 4\\ -4 & 0 \end{pmatrix}$, que es $A^2$ calculado
directamente.

\textbf{10.} Inducción:
$F^1 = \begin{pmatrix} F_2 & F_1\\ F_1 & F_0\end{pmatrix}$, y
\[
F^{n+1} = F^n F =
\begin{pmatrix} F_{n+1} + F_n & F_{n+1}\\
F_n + F_{n-1} & F_n \end{pmatrix}
= \begin{pmatrix} F_{n+2} & F_{n+1}\\ F_{n+1} & F_n
\end{pmatrix}.
\]
Aquí $s = 1$, $p = -1$, $D = X^2 - X - 1$ con raíces $\varphi,
\psi$ ($\varphi - \psi = \sqrt5$, $\varphi\psi = -1$). La sucesión
$(F_n)$ tiene $F_0 = 0 = a_0$, $F_1 = 1 = a_1$ y obedece la misma
recurrencia que $(a_n)$: $F_n = a_n = (\varphi^n - \psi^n)/\sqrt5$,
la fórmula de Binet. Cassini: aplicando la multiplicatividad de la
pregunta 7 a $F^n$,
\[
F_{n+1}F_{n-1} - F_n^2 = p(F^n) = p(F)^n = (-1)^n .
\]

\textbf{11.} La condición es \omterm{def:b1:linmaps:def}{lineal} y contiene la sucesión nula: un
\omterm{def:b1:vspaces:subspace}{subespacio}. Por inducción, $u_0, u_1$ determinan $u$ de \omterm{def:b1:linmaps:forms}{forma
lineal}, y cada pareja de valores iniciales la realiza exactamente una
solución: como en el~\cref{exo:b1:findim:10}, $E_D$ está
parametrizado \omterm{def:b1:logic:inj}{biyectiva} y linealmente por $(u_0, u_1) \in K^2$:
$\dim E_D = 2$.

\textbf{12.} $(a_n)$ obedece la recurrencia (pregunta 2) con
$a_0 = 0$, $a_1 = 1$. Y $(b_n)$ también:
$b_{n+2} = -p\,a_{n+1} = -p(s a_n + b_n) = s\,b_{n+1} - p\,b_n$
(usando dos veces $b_{n+1} = -pa_n$), con $b_0 = 1$, $b_1 = 0$. La
combinación $v_n = u_1 a_n + u_0 b_n$ es entonces una solución con
$v_0 = u_0$, $v_1 = u_1$; y dos soluciones con los mismos valores
iniciales coinciden (inducción), luego $u_n = u_1 a_n + u_0 b_n$ para
todo $n$.

\textbf{13.} $(\lambda^n)$ es solución si y solo si
$\lambda^{n+2} = s\lambda^{n+1} - p\lambda^n$ para todo $n$, es
decir, $D(\lambda) = 0$ (tras dividir por $\lambda^n \neq 0$;
nótese que $\lambda, \mu \neq 0$, ya que $p = \lambda\mu \neq 0$).
Libertad de $\bigl((\lambda^n), (\mu^n)\bigr)$: una relación en
$n = 0, 1$ da $c + c' = 0$, $c\lambda + c'\mu = 0$, luego
$c(\lambda - \mu) = 0$: $c = c' = 0$. Dos vectores \omterm{def:b1:vspaces:free}{libres} en
dimensión $2$: una \omterm{def:b1:vspaces:free}{base}.
Raíz doble: $\bigl((n\lambda^n)\bigr)$ es solución, ya que, con
$s = 2\lambda$, $p = \lambda^2$:
\[
s(n+1)\lambda^{n+1} - p\,n\lambda^n
= \lambda^{n+2}\bigl(2(n+1) - n\bigr) = (n+2)\lambda^{n+2} ;
\]
libertad en $n = 0, 1$: $c = 0$ y después $c'\lambda = 0$ con
$\lambda \neq 0$.

\textbf{14.} $D = X^2 - X - 6 = (X - 3)(X + 2)$. Solución general
$u_n = A\,3^n + B(-2)^n$; las condiciones iniciales dan $A + B = 1$ y
$3A - 2B = 8$, luego $A = 2$, $B = -1$:
\[
u_n = 2\cdot 3^n - (-2)^n .
\]
Comprobación: $u_2 = 18 - 4 = 14 = u_1 + 6u_0$;
$u_3 = 54 + 8 = 62 = u_2 + 6u_1 = 14 + 48$.

\textbf{15.}
$C\begin{pmatrix} u_n\\ u_{n+1}\end{pmatrix} =
\begin{pmatrix} u_{n+1}\\ -p\,u_n + s\,u_{n+1}\end{pmatrix} =
\begin{pmatrix} u_{n+1}\\ u_{n+2}\end{pmatrix}$, y la inducción da la
fórmula con $C^n$. Además, $\operatorname{tr} C = 0 + s = s$ y
$p(C) = 0\cdot s - 1\cdot(-p) = p$: la matriz compañera tiene
exactamente $D$ como \omterm{def:b1:poly:def}{polinomio} de Cayley--Hamilton.

\textbf{16.} Para cualquier solución $u$ de
$u_{n+3} = \alpha u_{n+2} + \beta u_{n+1} + \gamma u_n$, los vectores
de estado $v_n = (u_n, u_{n+1}, u_{n+2})^{\mathsf T}$ cumplen
$C_3 v_n = v_{n+1}$ (las dos primeras filas desplazan y la última
aplica la recurrencia). Por tanto,
\[
D_3(C_3)\,v_0 = v_3 - \alpha v_2 - \beta v_1 - \gamma v_0 ,
\]
cuyas tres componentes son
$u_{k+3} - \alpha u_{k+2} - \beta u_{k+1} - \gamma u_k = 0$
($k = 0, 1, 2$). Y como el estado inicial
$v_0 = (u_0, u_1, u_2)^{\mathsf T}$ recorre \emph{todo} $K^3$ (los
valores iniciales son \omterm{def:b1:vspaces:free}{libres}), la matriz $D_3(C_3)$ mata a todos los
vectores: $D_3(C_3) = 0$.

\textbf{17.} Escríbase $X^n = Q\,D_3 + R_n$ con $\deg R_n \leq 2$ y
evalúese en cada raíz: $\lambda_i^n = R_n(\lambda_i)$. Así, $R_n$ es
un \omterm{def:b1:poly:def}{polinomio} de grado $\leq 2$ que interpola los tres valores
$\lambda_i^n$ en los tres nodos distintos $\lambda_i$: por la
unicidad del~\cref{thm:b1:poly:lagrange},
$R_n = \sum_i \lambda_i^n L_i$, con $(L_i)$ la \omterm{def:b1:vspaces:free}{base} de Lagrange de
los nodos. Sustituyendo $C_3$ (preguntas 5 y 16):
\[
C_3^{\,n} = R_n(C_3) = \sum_{i=1}^{3} \lambda_i^n\,L_i(C_3),
\]
con las tres matrices $L_i(C_3)$ independientes de $n$: toda entrada
de $C_3^{\,n}$ es una combinación fija de
$\lambda_1^n, \lambda_2^n, \lambda_3^n$.

\textbf{18.} $D_3 = X^3 - 2X^2 - X + 2 = (X-1)(X+1)(X-2)$. Solución
general $u_n = A + B(-1)^n + C\,2^n$. Condiciones iniciales:
$A + B + C = 0$, $A - B + 2C = 1$, $A + B + 4C = 1$. Restando la
primera de la tercera: $3C = 1$, $C = \frac13$; entonces
$A + B = -\frac13$ y $A - B = \frac13$: $A = 0$, $B = -\frac13$. Por
tanto,
\[
u_n = \frac{2^n - (-1)^n}{3}
\]
(los números de Jacobsthal). Comprobación:
$u_3 = \frac{8 + 1}{3} = 3 = 2u_2 + u_1 - 2u_0 = 2 + 1 - 0$.

\textbf{19.} Desarrollo de Taylor del \omterm{def:b1:poly:def}{polinomio} $X^n$ en
$\lambda$:
\[
X^n = \sum_{k=0}^{n} \binom nk \lambda^{n-k}(X - \lambda)^k ,
\]
y todos los términos con $k \geq 3$ son divisibles por
$(X - \lambda)^3$: el resto es
\[
R_n = \lambda^n + n\lambda^{n-1}(X - \lambda) + \binom
n2\lambda^{n-2}(X - \lambda)^2 .
\]
Para $M = \lambda I + N$ con $N^3 = 0$:
$(M - \lambda I)^3 = N^3 = 0$, así que la pregunta 5 da
\[
M^n = \lambda^n I + n\lambda^{n-1} N + \binom n2
\lambda^{n-2} N^2 ,
\]
que es exactamente el desarrollo binomial de $(\lambda I + N)^n$
truncado en $N^2$ --- los dos métodos coinciden.

\textbf{20.} El espacio de soluciones tiene dimensión $3$ (la misma
parametrización por $(u_0, u_1, u_2)$ que en la pregunta 11), y cada
$(\lambda_i^n)$ es solución. Libertad: supóngase
$c_1\lambda_1^n + c_2\lambda_2^n + c_3\lambda_3^n = 0$ para
$n = 0, 1, 2$. Fíjese $i$ y sea $L_i = \sum_{k \leq 2} p_k X^k$ el
\omterm{def:b1:poly:def}{polinomio} de Lagrange de los nodos con
$L_i(\lambda_j) = \delta_{ij}$. Entonces
\[
0 = \sum_{k=0}^{2} p_k\Bigl(\sum_j c_j\lambda_j^k\Bigr)
= \sum_j c_j\,L_i(\lambda_j) = c_i .
\]
Así, todos los $c_i = 0$: tres soluciones \omterm{def:b1:vspaces:free}{libres} en dimensión $3$,
una \omterm{def:b1:vspaces:free}{base}; y la solución general es
$c_1\lambda_1^n + c_2\lambda_2^n + c_3\lambda_3^n$.

\textbf{21.} De $F_k = F_{k+2} - F_{k+1}$, la suma telescopa:
\[
\sum_{k=1}^{n} F_k = \sum_{k=1}^{n}\bigl(F_{k+2} - F_{k+1}\bigr)
= F_{n+2} - F_2 = F_{n+2} - 1 .
\]

\textbf{22.} Tómese la entrada $(1,2)$ de $F^{m+n} = F^m F^n$: el
miembro izquierdo es $F_{m+n}$; y el derecho es (la fila $1$ de
$F^m$) por (la columna $2$ de $F^n$), es decir,
$F_{m+1}F_n + F_m F_{n-1}$. Con $m = n$:
\[
F_{2n} = F_{n+1}F_n + F_nF_{n-1} = F_n\,(F_{n+1} + F_{n-1}).
\]

\textbf{23.} Por Binet,
$F_n - \dfrac{\varphi^n}{\sqrt5} = -\dfrac{\psi^n}{\sqrt5}$, y
$\abs\psi = \frac{\sqrt5 - 1}2 < 1$, luego
\[
\Bigl|F_n - \frac{\varphi^n}{\sqrt5}\Bigr|
\leq \frac{1}{\sqrt5} < \frac12
\qquad (n \geq 0):
\]
$F_n$ es el entero más próximo a $\varphi^n/\sqrt5$.

\textbf{24.} $t_n = F_{n+1} + F_{n-1}$ es una combinación de
sucesiones de Fibonacci desplazadas, luego cumple la misma
recurrencia: $t_{n+2} = t_{n+1} + t_n$; y $t_1 = F_2 + F_0 = 1$,
$t_2 = F_3 + F_1 = 3$: son los números de Lucas $L_n$. La sucesión
$\varphi^n + \psi^n$ es una solución con los mismos dos primeros
valores ($\varphi + \psi = 1$,
$\varphi^2 + \psi^2 = ( \varphi + \psi)^2 - 2\varphi\psi = 3$), luego
$L_n = \varphi^n + \psi^n$. Por último,
\[
F_n L_n = \frac{(\varphi^n - \psi^n)(\varphi^n +
\psi^n)}{\sqrt5} = \frac{\varphi^{2n} - \psi^{2n}}{\sqrt5} =
F_{2n},
\]
lo que recupera la pregunta 22.

\textbf{25.} (i) Las cinco matrices $I, A, A^2, A^3, A^4$ viven en el
espacio de dimensión $4$, $\mathcal{M}_2(K)$, así que \emph{algún}
\omterm{def:b1:poly:def}{polinomio} no nulo de grado $\leq 4$ mata a $A$; la parte II afinó
esto hasta la cuadrática explícita $A^2 = sA - pI$, que encierra
todas las potencias en el plano $\operatorname{Vect}(I, A)$. (ii) La
división euclídea reduce $X^n$ \omterm{def:b1:complex:field}{módulo} esa cuadrática, y los dos
coeficientes del resto obedecen la recurrencia de dos términos
$a_{n+2} = s\,a_{n+1} - p\,a_n$: la exponenciación se ha convertido
en iteración. (iii) La pregunta 6 es el caso $n = 2$ del
\emph{teorema de Cayley--Hamilton}, válido en toda dimensión y
demostrado en el volumen del segundo año. (iv) La matriz compañera
cierra el círculo: toda \omterm{pb:b1:matrices:1}{recurrencia lineal} \emph{es} una potencia de
matriz, con el mismo \omterm{def:b1:poly:def}{polinomio} $D$ apareciendo como datos de \omterm{def:b1:matrices:transpose}{traza} y
determinante, de modo que el cálculo con restos resuelve
recurrencias y calcula potencias de un solo golpe.
\end{solution}
